\section{Metropolis}
\subsection{Simulation Workflow: From Batch to Single Run}

Two functions are used for running the simulation: a high-level manager 
(\texttt{mega\_metropolis}), and a core simulation engine (\texttt{run\_metropolis}). 
This design allows for parallelized execution of simulations on multiple 
initial murine sequences.

\subsubsection{The \texttt{mega\_metropolis} Manager}

\texttt{mega\_metropolis}'s primary responsibilities are:

\begin{enumerate}
    \item \textbf{Data Loading:} It begins by reading a list of initial murine sequences 
        (which we call seeds) from a specified file and saves them for later use.
    \item \textbf{Batch Management:} It creates a timestamped main directory for the entire 
        simulation batch (format \texttt{results/metropolis/YYYY-MM-DD\_HH\_MM\_SS\_.../}). .
    \item \textbf{Parallel Execution:} It iterates through each seed sequence using an OpenMP parallel loop. 
    For each seed, it creates a dedicated subdirectory (e.g., \texttt{seq\_1}, \texttt{seq\_2}).
    \item \textbf{Replication:} For each seed, it calls the \texttt{run\_metropolis} function \texttt{n\_metropolis} times. 
        This allows for multiple independent simulation runs starting from the same seed, providing statistical 
        significance to the results. The output of each run is saved to a separate file (e.g., \texttt{run\_1.txt}, 
        \texttt{run\_2.txt}) within the corresponding seed's subdirectory.
\end{enumerate}

\subsubsection{The \texttt{run\_metropolis} Engine}

The \texttt{run\_metropolis} function performs a single, complete simulated annealing run on one given sequence:

\begin{enumerate}
    \item \textbf{Initialization:} It takes an initial sequence, the human Schrödinger chain, and the 
        pre-calculated matrix of $\beta$ values.
    \item \textbf{Simulated Annealing Loop:} It iterates through each row of the $\beta$ matrix, 
        representing a step in the ``cooling'' schedule.
    \item \textbf{Metropolis Sweeps:} For each $\beta^i$ step, it performs a specified number of 
        Metropolis sweeps (\texttt{n\_sweeps}). A single sweep involves attempting a mutation at every position in the sequence.
    \item \textbf{State Acceptance:} We begin by proposing a new state, selecting a new AA (including the gap character) 
        uniformly from the 21 possibilities (ensuring it's different from the current one). 
        Each proposed mutation is accepted or rejected based on the change in total energy 
        ($\Delta E$) and the current, position-specific, $\beta^i_j$ value, following the Metropolis criterion: 
        accept if $\Delta E < 0$ or if a random number is less than $e^{-\beta^i_j \Delta E}$.
    \item \textbf{Data Logging:} After the simulation is complete, it calls a dedicated function to write a detailed report 
        file containing the initial parameters, the final sequence and energy for each $\beta$ step, 
        and overall acceptance statistics.
\end{enumerate}

\subsection{Simulated Annealing Schedule ($\beta$)}
The $\beta$-values, which are inverse of the temperature ($1/k_B T$), are not constant; instead, they are defined through 
% two different mechanisms (the ``engine'' and the ``damper'').
a mechanism, which we call the ``engine'':
\begin{itemize}
    \item \textbf{Engine:} 
    $$\displaystyle\beta_i^j = k_i \cdot \frac{\alpha}{S_j + \varepsilon}, \quad k_i = \left(c_r\right)^i\ 
    \forall i \in [1,N_\beta], j \in [1, L_\text{chain}]$$ 
    where $c_r$ is the cooling rate, $\varepsilon$ avoids division by 0, $\alpha$ is a scale factor, and $S_j$ is 
    entropy for position $j$ in human reference Schrödinger chain.
    % \item \textbf{Damper:} 
    % $$\varphi(\sigma, \lambda) = \frac{L}{\displaystyle 1+(1-L)e^{-\lambda \sigma}}$$ 
    % where $\lambda$ is a global variable (the dampening coefficient), $L$ is a global variable too (maximum value our $\beta$ 
    % will be multiplied by), and $\sigma$ is the deviation of the last sweep acceptance from the target acceptance. 
    % Since we want acceptance to decay as our system explores the energy landscape, we begin with a \texttt{STARTING\_TARGET\_ACCEPTANCE}
    % value, $\alpha_0$, (defined in \texttt{head.h}), and slowly decrease it using the following equation: 
    % $$\alpha_i = \alpha_0 \left(1-\frac{i}{N_\beta}\right)$$ 
    % where $N_\beta$ is the number of betas we have, and $i$ is the current iteration 
    % with our pre-defined $\beta$ given by the engine.
    % $\newline$ Lastly, $\sigma = \alpha_{ij}-\alpha_i$, where $\alpha_{ij}$ is 
    % acceptance from last sweep iteration $(j\in[1,\text{n\_sweeps}])$
\end{itemize}

\subsection{Energy Function}
The total energy, $E_\text{total}$, is a weighted sum of three components, which we aim to minimize:
    $$E_\text{total} = w_\text{log} E_\text{log} + w_\text{prop} E_\text{prop} + w_\text{penalty} E_\text{penalty}$$
    The weights for each component (\texttt{WEIGHT\_LOG}, \texttt{WEIGHT\_PROP}, \texttt{WEIGHT\_PENALTY}) are defined in `head.h`.
    \begin{enumerate}
        \item \textbf{Log Humanness Energy, }$\mathbf{E_\text{log}}$:
        This term quantifies how "human-like" a sequence is based on AA frequencies from a 
        reference set of human sequences (the Schrödinger chain):
        $$E_\text{log} = -\sum_{i=1}^{L_\text{chain}} \log(p_i(a_i))$$
        where $p_i(a_i)$ is the frequency of the amino acid $a_i$ at position $i$ in the human reference. 
        A large penalty (\texttt{ZERO\_FREQ\_PENALTY\_LOG}) is added for amino acids with a frequency below a 
        small threshold (\texttt{EPSILON}).
        \item \textbf{Property Distance Energy, }$\mathbf{E_\text{prop}}$:
        This term measures the similarity between the properties of the sequence and the human reference; 
        it's calculated as the squared Euclidean distance in the property space for each position:
        $$E_\text{prop} = \sum_{i=1}^{L_\text{chain}} \sum_{j=1}^{N_\text{props}} \left( \text{prop}_j(a_i) - 
        \overline{\text{prop}_j(i)} \right)^2$$
        where $\text{prop}_j(a_i)$ is the value of property $j$ for amino acid $a_i$, and $\overline{\text{prop}_j(i)}$ 
        is the average value of property $j$ at position $i$ in the Schrödinger chain.
        \item \textbf{Wandering Penatly Energy, }$\mathbf{E_\text{penalty}}$:
        This term discourages the sequence from deviating too far from the original murine sequence; 
        we thus only incentivize mutation if there's a relatively big gain in humanness. 
        It consists of a penalty based on property distance and a (for now) constant penalty for any mutation:
        $$E_\text{penalty} = \sum_{i=1}^{L_\text{chain}} \left[ \left( \sum_{j=1}^{N_\text{props}} 
        (\text{prop}_j(a_i) - \text{prop}_j(a_i^\text{original}))^2 \right) + P \cdot 
        \delta_{\displaystyle  a_i}^{\displaystyle a_i^{\text{original}}}\right]$$
        where $a_i^\text{original}$ is the amino acid at position $i$ in the original murine sequence, 
        $P$ is the \texttt{AA\_MUTATION\_PENALTY}, and $\delta$ is the ``generalized Kronecker delta'' (logical if gate). 
        This approach ensures each mutation has a minimum energy cost, while heavily penalizing big property changes.
    \end{enumerate}

\subsection{Model Parameters}
The following table summarizes the key parameters for metropolis in our simulations. These values are defined in \texttt{main.c} and \texttt{head.h}.

\begin{table}[h!]
\centering
\begin{tabular}{|l|l|p{7cm}|}
\hline
\textbf{Parameter} & \textbf{Value (provisional)} & \textbf{Description} \\
\hline
\texttt{n\_betas} & 50 & Number of cooling steps in the annealing schedule. \\
\texttt{n\_sweeps} & 800 & Number of Metropolis sweeps per $\beta$ step. \\
\texttt{n\_metropolis} & 20 & Number of independent runs per murine seed sequence. \\
\texttt{w\_log} & 0.4 (from head.h) & Weight for the log-humanness energy component. \\
\texttt{w\_prop} & 0.4 (from head.h) & Weight for the property-distance energy component. \\
\texttt{w\_penalty} & 0.2 (from head.h) & Weight for the wandering penalty energy component. \\
\texttt{cooling\_rate ($c_r$)} & 1.02 & Multiplicative factor for increasing $k_i$ at each step. \\
\texttt{scale\_factor ($\alpha$)} & 1.0 & Global scaling factor for all $\beta$ values. \\
\texttt{AA\_MUTATION\_PENALTY ($P$)} & 1.5 (from head.h) & Constant penalty added for any AA mutation away from the original. \\
\hline
\end{tabular}
\label{tab:metropolis_params}
\end{table}